% =============================================================================
% AGENTE  : Cientista-Líder + Geneticista (Especialista em Domínio)
% MÓDULO  : Introdução — versão em Português Brasileiro (intro_pt.tex)
% SINCRON.: Deve ser consistente com intro_en.tex (responsabilidade do Cientista-Líder)
% STATUS  : RASCUNHO
% =============================================================================

\section{Introdução}
\label{sec:intro}

A caracterização da diversidade genética constitui um pilar central do
melhoramento de plantas, da genética da conservação e da biologia
evolutiva \autocite{Hamrick1989}. Os marcadores moleculares tornaram-se
ferramentas indispensáveis para a avaliação da diversidade em nível
genômico, permitindo a identificação de polimorfismos com elevada
precisão, reprodutibilidade e capacidade analítica
\autocite{Semagn2006}.

Dentre as tecnologias de marcadores atualmente disponíveis, os
Polimorfismos de Nucleotídeo Único (\snp{}) e as Repetições de
Sequências Simples (\ssr{}) são as mais amplamente empregadas. Os
marcadores \snp{} representam a classe mais abundante de variação
genética em qualquer genoma e são compatíveis com plataformas de
genotipagem de alto rendimento, como \textit{arrays} de \snp{} e
genotipagem por sequenciamento (GBS, do inglês
\textit{Genotyping-by-Sequencing}) \autocite{Elshire2011}. Os marcadores
\ssr{}, também conhecidos como microssatélites, são altamente
polimórficos, codominantes e amplamente validados em numerosas espécies
cultivadas, sendo valiosos para mapeamento genético, análise de
parentesco e identificação de variedades \autocite{Powell1996}.

A informatividade de um loco marcador é comumente quantificada pelo
Conteúdo de Informação sobre Polimorfismo (\pic{}), proposto
originalmente por \textcite{Botstein1980}. O valor de \pic{} varia de 0
(monomórfico) a um máximo teórico de aproximadamente 0,5 para marcadores
bialélicos, sendo valores acima de 0,5 geralmente considerados
altamente informativos para locos multialélicos. A heterozigose esperada
(\he{}) e a observada (\ho{}) são estatísticas complementares que
refletem a distribuição das frequências alélicas dentro de uma população
e os desvios do equilíbrio de Hardy--Weinberg, respectivamente.

Apesar da ampla aplicação de ambas as classes de marcadores na
literatura, análises comparativas de seu poder discriminatório sob
condições de simulação padronizadas permanecem pouco exploradas.
Compreender o desempenho relativo dos marcadores \snp{} e \ssr{} em
termos de \pic{} e heterozigose é essencial para o delineamento de
estratégias de genotipagem eficientes em estudos futuros.

O presente estudo tem como objetivos: (i)~simular um painel
representativo de marcadores \snp{} e \ssr{} com distribuições de
frequência alélica biologicamente plausíveis; (ii)~calcular \pic{},
\he{} e \ho{} para cada classe de marcador; e (iii)~fornecer uma
comparação visual e estatística para embasar critérios de seleção de
marcadores em programas de melhoramento aplicado.
